\section{Proyección y escalamiento}
\label{sec:proyeccion_escalamiento}

La modularidad MVC y la selección de herramientas de código abierto (Vista HTML5/CSS3/JS, Controlador PHP 8.2/Apache, Modelo PDO--PostgreSQL 16.4) otorgan al sistema una elevada capacidad de réplica, portabilidad y escalamiento geográfico e institucional:

\subsection{Escalamiento geográfico y transferencia tecnológica}

Al no depender de licencias propietarias ni de plataformas de pago recurrentes, el software desarrollado puede ser transferido sin costo alguno a otros Módulos Tecnológicos Productivos y unidades educativas técnico-humanísticas del departamento de Santa Cruz (tales como los centros BTH de los municipios vecinos de Cuatro Cañadas, San Ramón, San Javier y Guarayos). Estas instituciones comparten una estructura sociocomunitaria y desafíos organizativos similares durante sus procesos anuales de matrícula escolar, beneficiándose de un modelo validado en el terreno.

\subsection{Escalabilidad de infraestructura tecnológica}

En el plano de la infraestructura, se proyecta la migración progresiva desde el servidor web local actual hacia un entorno de alojamiento en la nube sobre un Servidor Virtual Privado (VPS) con enlace dedicado, respaldos automatizados programados mediante utilitarios de PostgreSQL (pg\_dump) y certificados de seguridad SSL/TLS para el cifrado extremo a extremo de las comunicaciones. La estimación referencial asciende a 800 Bs por año por VPS (1 vCPU/1 GB RAM) más SSL gratuito vía Let's Encrypt, sin costos de licencias propietarias; frente al ahorro colectivo superior a 20000 Bs por gestión y a la inversión inicial de 13556 Bs, el retorno se mantiene en menos de una gestión incluso en réplica multi-sede.

De manera complementaria, se proyectan como nuevas funcionalidades técnicas de la Versión 3.0 las notificaciones automatizadas por WhatsApp para alertas de credenciales y estado de expediente, el módulo de pre-auditoría documental con carga preliminar de fotografías/escaneos de libretas de 1ro a 3ro, cédula de tutores, carnet y certificado, y la emisión de credenciales QR para control de acceso a talleres. Estas mejoras, previstas en la ficha del proyecto, responden de forma directa a la heterogeneidad de competencias digitales evidenciada en el 8\% de usabilidad pendiente y a la necesidad de reducir la contrastación presencial en secretaría.

\subsection{Alianzas estratégicas y gobernanza distrital}

Para garantizar la viabilidad a largo plazo de la plataforma, se propone la formalización de un convenio interinstitucional entre el MÓDULO TECNOLÓGICO PRODUCTIVO SAN JULIÁN, la Dirección Distrital de Educación de San Julián y el Gobierno Autónomo Municipal de San Julián. Dicha alianza facilitará la sincronización de los calendarios escolares, la provisión de conectividad permanente en los predios del módulo y la articulación con las direcciones de las unidades educativas humanísticas de origen para la validación cruzada y oportuna de los estudiantes matriculados en cada curso y paralelo.
